\documentclass[12pt]{article}

\usepackage[margin=1in]{geometry}
\usepackage{setspace}
\usepackage[numbers,sort&compress]{natbib}
\usepackage{amsmath, amssymb, amsthm}
\usepackage{times}
\usepackage[hidelinks]{hyperref}
\usepackage{float}
\usepackage{gensymb}

\onehalfspacing

\title{Heat-related mortality risk across Canada's
neighbourhoods: a case time series analysis of 41 census
metropolitan areas}
\author{[Author line]}
\date{}

\begin{document}
\maketitle

% LOG --- revise each session
%============================================================
%
% STATUS: full end-to-end draft, awaiting next week pass
% ---- LAST SESSION (2026-07-13) ----
% - Drafted the whole paper.
% - Title is a placeholder---re-derive from most surprising result.
%
% ---- OPEN FOLLOW-UPS ----
% [ ] refs.bib --- Zotero + Better BibTeX export. Citekeys locked in-text.
%     NEW citekeys from proposal, confirm in .bib: hebbern2023, cote2024,
%     and possibly rinner2010 (Toronto, uncited in P3 --- cite or leave).
% [ ] Add \usepackage{gensymb} to preamble --- \degree won't compile.
% [ ] [FILL] numbers: Boudreault Quebec figure; Fig 1/2/4 RR + CI +
%     ranges; Fig 2 Q p-value + I2; Fig 4 RR range + spatial pattern.
% [ ] Fig 3 recovery: re-run after DGP collinearity fix (F1/F3 too close).
% [ ] Author line + affiliations.
% [ ] Fig 3 slot changes job with real mortality (no planted truth):
%     decide descriptive observed PCA vs validation-to-appendix.
%
% ---- ACKNOWLEDGMENTS (important, do not lose) ----
% CANUE attribution required by the data-use agreement. Check their
% metadata for the required wording and how to cite them. Put the
% CANUE statement in the Acknowledgments at the end.
%
% ---- FIGURE SLOT INVARIANT (pilot -> final) ----
% Slot 1 fits: slivers -> full/national Stage 1
% Slot 2 pool: 3 CMAs / age 75-84 -> all CMAs, all age bands
% Slot 3 recovery: simulated -> observed structure OR appendix
% Slot 4 map: pilot geography -> national DA map
%

%============================================================
% ABSTRACT --- written last; four parts, eight sentences.
%============================================================
\begin{abstract}
%Background:
Extreme heat kills in Canada far beyond the deaths certified as
heat, and its risk concentrates in particular neighbourhoods. A neighbourhood
records too few deaths to estimate its risk directly, and
national estimates have therefore resolved to the city and
above.
%Methods:
We applied a two-stage small-area design to the [XX]
dissemination areas of Canada's 41 census metropolitan areas.
Case time series fits within each city and age band were pooled
in a multivariate meta-regression on composite vulnerability
components and downscaled to an exposure-response curve for
every dissemination area.
%Results:
Heat-related mortality risk varied [more within cities than
between them], with relative risk at the 99th temperature
percentile ranging from [XX] to [XX] across dissemination areas
and [XX] excess deaths (95\% empirical CI [XX--XX]) attributable
to heat each year. [Vulnerability components / the dominant
modifier] accounted for [XX]\% of the between-area variation in
heat risk.
%Conclusion:
The analysis yields a national map of heat-related mortality
risk at neighbourhood resolution, with minimum-mortality
temperature and attributable burden for every dissemination
area. We recommend the dissemination-area estimates as the
input layer for municipal heat-response planning.
% [GATE: Results S1 is the surprise slot (criterion 1, may
%  flip). S2 filled from Fig 3. DA count ~40,000 in-CMA, pin at
%  extraction. Eight sentences, at ceiling.]
\end{abstract}

%============================================================
\section{Introduction} % - field, ladder, routes, aim.
%============================================================

Non-optimum ambient temperature accounts for 7$\cdot$71\% of
mortality across 384 locations in 13 countries, Canada among
them \citep{gasparrini2015mcc}, and more than a third of
warm-season heat-related deaths worldwide are attributable to
human-induced warming \citep{vicedocabrera2021}. In Canada,
roughly 670 excess deaths over 2000 to 2020 have been attributed
to extreme heat events across the country's twelve largest
cities \citep{quick2024twelvecities}. The June 2021 heat dome
was associated with [619] deaths in British Columbia
\needcite{BC Coroners heat dome death review---verify count
and dates against presentation}. Province-wide assessments place
a burden of comparable scale within single provinces
\citep{boudreault2024quebec}. Heat risk is not spread evenly–it
concentrates in particular places and populations, and
area-level characteristics account for much of the difference.
Across Europe, small-area studies find heat mortality higher
where population density is greater, green space scarcer, and
air pollution higher \citep{zhang2023cvdeurope,
zafeiratou2023respeurope}. The same characteristics recur across
cities and countries, with only their local weights differing
\citep{sera2019urban}.

Canadian heat-mortality work has operated above the
neighbourhood scale. National modelling has estimated
temperature--mortality associations across 111 health regions
\citep{hebbern2023}, multi-city work across the twelve largest
cities \citep{quick2024twelvecities}, and provincial analyses
within single provinces \citep{boudreault2024quebec,
chen2016ontario}. Neighbourhood-scale work has been single-city:
Toronto mapped through composite vulnerability indices
\citep{rinner2010}, Quebec's dissemination areas through
aggregate counts \citep{cote2024}, and three cities through the
national death registry \citep{quick2025smallarea}.

Estimation below the city is hard because the unit records too
few deaths to estimate a stable exposure-response curve of its
own. Two routes get past it. Single-stage spatial models fit all
areas at once, borrowing strength between neighbours through a
conditional autoregressive prior \citep{quijalzamorano2024sbdlnm,
rutten2026spatialvarying, rutten2026effectmod}. Two-stage
designs fit where counts allow, pool the reduced curves in a
multivariate meta-regression on area-level vulnerability
predictors, and predict every area from its own profile
\citep{gasparrini2022smallarea, sera2022extended,
dimaso2025italy, masselot2025extensions}. A smoothing prior
pools neighbours hardest in the temperature tail, where
vulnerability binds and heat risk is read; the two-stage design
carries no geographic term, so each area's risk follows its own
profile.

We apply the two-stage small-area design
\citep{gasparrini2022smallarea} to the [XX] dissemination areas
of Canada's 41 census metropolitan areas: a case time series fit
within each city and age band, pooled across cities on composite
vulnerability components, and downscaled to an exposure-response
curve for every dissemination area. Doing so asks two questions
of the Canadian substrate: to what extent heat-related mortality
risk varies across neighbourhoods, and how much of that
variation the vulnerability indicators account for.

%============================================================
\section{Methods} % - fit → reduce → pool → downscale.
%============================================================

\subsection{Study design and data}

We performed a nationwide small-area analysis of mortality and
temperature across the [XX] dissemination areas (DAs) within
Canada's 41 census metropolitan areas (CMAs), from 2000 to 2023.

Individual death records from the Canadian Vital Statistics
Death database (CVSD), obtained through Statistics Canada's
virtual Data Lab, were assigned to DAs by residential postal
code through the PCCF+ linkage and aggregated to daily counts of
all-cause non-accidental deaths in four age bands (0--64, 65--74,
75--84, 85 and over). Each DA carries a daily mean temperature
series, the mean of daily minimum and maximum from the Daymet V4
1 km gridded product, area-weighted over its polygon.
Age-specific population counts come from the 2021 Census, on
2021 boundaries throughout. DAs with zero population or without
temperature coverage were excluded.

Area-level vulnerability indicators come from thirteen datasets
of the Canadian Urban Environmental Health Research Consortium
(CANUE; greenness, air quality, urban heat, built form, 2019),
indexed to single-link postal codes at a 100 m buffer and
assigned to DAs through the same PCCF+ linkage, and from
thirteen topics of the 2021 Census Profile (demography, housing,
deprivation), taken as rates. Air conditioning is not collected;
renter share stands in for it. The Pampalon Material and Social
Deprivation Index benchmarks the components and is not an input,
as it is built from census variables that enter directly.

\subsection{Statistical analysis}

Daily death and temperature series enter as a case time series
per DA \citep{gasparrini2021, gasparrini2022cts}, fitted
separately for each CMA and age band, giving 164 first-stage
estimates.

Each fit is a conditional Poisson model with quasi-Poisson
variance, with stratum intercepts eliminated by conditioning
\citep{armstrong2014}. Strata are defined by DA, year, and
month, with day-of-week indicators and daily relative humidity
[source] as covariates and no further smooth of time. Temperature enters through a cross-basis
\citep{gasparrini2014dlnm}: a quadratic B-spline with three
internal knots at the 10th, 75th, and 90th percentiles of the
CMA's temperature distribution, and a natural cubic spline over
lags 0 to 21 with three internal knots equally spaced on the log
scale. The two margins give five basis columns each, and the
cross-basis 25 parameters.

We reduced each fitted cross-basis to its overall cumulative
exposure-response curve, centred at the CMA's median
temperature, yielding a five-coefficient summary and its
covariance, the input to the second stage
\citep{gasparrini2013reducing}. First-stage fits and reductions
ran inside the virtual Data Lab; only the reduced coefficients
and their covariances left it.

Vulnerability enters through a single principal component
analysis fitted on the DA rows of all 41 CMAs; the first three
components are carried forward, and CMA-level scores are
population-weighted averages of the DA scores.

We pooled the 164 reduced curves in a multivariate
meta-regression under restricted maximum likelihood, with age
band indicators and the first three components as
meta-predictors and a random intercept on CMA
\citep{sera2019mixmeta, gasparrini2012mvmeta, sera2022extended}.
Residual heterogeneity is summarized by Cochran's Q and
I\textsuperscript{2}.

From the pooled model we predicted a curve for every DA and age
band from the DA's own component scores, in two passes: a first
prediction at the CMA-median centring locates the DA's
minimum-mortality temperature, and the curve re-centred there
yields the relative risks of heat and cold at the 99th and 1st
percentiles and the corresponding attributable fractions
\citep{gasparrini2014attributable}. Excess mortality
rates are standardized to the Canadian 2011 standard population.
Empirical confidence intervals come from 1,000 Monte Carlo draws
of the pooled coefficients from their estimated covariance, with
the per-DA prediction repeated on each draw, CMA and national
totals summed within each draw, and intervals taken at the 2.5th
and 97.5th quantiles \citep{masselot2025extensions}.

Sensitivity analyses varied the lag window (0--14, 0--28 days),
the cross-basis dimension (9, 16), the exposure metric (daily
maximum; apparent temperature), the humidity adjustment
(excluded), the stratification (day of week absorbed into the
strata), the time control (a natural cubic spline of time
added), the season (warm season only), the CANUE buffer (postal
code point value), the standard population (European Standard
Population 2013), and the map spread (corrected for residual
between-area variance). between-area variance). Secondary analyses repeated the model by
cause of death and by sex. Analyses were performed in R
[version] with the dlnm, gnm, and mixmeta packages.
% [GATE: R version at final run.]
%============================================================
\section{Results} % how many died, who dies, why they differ, where they are.
%============================================================

The analysis included [XX] deaths from all causes across the 41
census metropolitan areas, spanning [XX] dissemination areas.
Each year over the study period, an average of [XX] excess deaths
(95\% empirical CI [XX--XX]) were attributable to heat,
corresponding to a standardized excess mortality rate of [XX]
deaths per 100,000 person-years (Canadian 2011 standard
population). Figure~\ref{fig:cities} shows the city-level heat
relative risk at the 99th temperature percentile across the 41
CMAs, ranging from [XX] in [CMA] to [XX] in [CMA].
% [GATE: CVSD extraction + national Stage 1 (DRAC). DA count:
%  ~40,000 within the 41 CMAs; ~57,900 is all-Canada -- pin the
%  frame before filling.]

% Figure 1 -- slot 1 endpoint (pilot: 5 sliver fits)
\begin{figure}[H]
\centering
% \includegraphics[width=\textwidth]{fig1_cma_rr.pdf}
\caption{City-level relative risk of heat at the 99th temperature
percentile across the 41 census metropolitan areas, from the
first-stage fits. [Placeholder---gated on observed mortality.]}
\label{fig:cities}
\end{figure}

Risk increased with age. The pooled relative risk at the 99th
percentile relative to the minimum-mortality temperature rose
from [XX] (95\% empirical CI [XX--XX]) at ages 0--64 to [XX]
([XX--XX]) at ages 85 and over (Figure~\ref{fig:pooled}).
Minimum-mortality temperatures ranged from [XX] to
[XX]$\,\degree$C across CMAs, tracking each city's warm-season
temperature distribution. Residual between-CMA heterogeneity
after meta-regression was [Q p-value], I\textsuperscript{2}
[XX]\%.
% [GATE: all-CMA, all-age Stage 2. Q/I2 computable at K=41
%  (df.residual > 0); the K=5 pilot pool cannot test them --
%  monograph carries that fact, not the paper.]

% Figure 2 -- slot 2 endpoint (pilot: K=5, one band)
\begin{figure}[H]
\centering
% \includegraphics[width=\textwidth]{fig2_pooled_age.pdf}
\caption{Pooled overall cumulative exposure-response curves by
age band, from the multivariate meta-regression across the 41
CMAs. [Placeholder---gated on observed mortality.]}
\label{fig:pooled}
\end{figure}

The first three components accounted for [XX]\% of the variance
of the indicators and entered the pooled model as
meta-predictors, with [which components] associated with the
heat slope; [component] tracked the Pampalon material
deprivation index (r = [XX]). Predicted curves for
a high-vulnerability and a low-vulnerability profile within the
same CMA differed by [XX]\% in relative risk at the 99th
percentile (Figure~\ref{fig:vuln}).
% [GATE: real-indicator two-level PCA + K=41 Stage 2. The >20%
%  within-CMA contrast is pilot criterion 1 promoted to a
%  reported estimate. Slot-3 job change executed: descriptive
%  observed structure, not recovery -- recovery lives in the
%  appendix validation note.]

% Figure 3 -- slot 3 endpoint (pilot: PCA recovery check)
\begin{figure}[H]
\centering
% \includegraphics[width=\textwidth]{fig3_vulnerability.pdf}
\caption{Vulnerability structure and its effect on risk:
principal component loadings of the area-level indicators, and
predicted exposure-response curves for high- and
low-vulnerability profiles within one CMA.
[Placeholder---gated on observed mortality.]}
\label{fig:vuln}
\end{figure}

Figure~\ref{fig:map} maps heat-related mortality risk at
dissemination-area resolution nationally. Predicted relative
risk at the 99th percentile ranged from [XX] to [XX] across DAs,
and minimum-mortality temperatures from [XX] to
[XX]$\,\degree$C. Variation within CMAs [exceeded / compared to]
variation between them, with the highest risks concentrated in
[spatial pattern]. Downscaled predictions agreed with unpooled
census-tract reference fits within [XX] (appendix).
% [GATE: national downscale on observed mortality. The
%  within/between spread sentence is evaluation criterion 1; the
%  census-tract sentence is the principle-9 transfer check --
%  appendix figure. Pilot Stage 3 diagnostics (Toronto 7,682 DAs,
%  |rho|~0.91 ordering, sign inversion open, 2.9x amplitude) are
%  monograph material and do not enter here.]

% Figure 4 -- slot 4 endpoint: THE money figure
\begin{figure}[H]
\centering
% \includegraphics[width=\textwidth]{fig4_national_map.pdf}
\caption{Heat-related mortality risk across Canadian
dissemination areas: predicted relative risk of heat at the 99th
temperature percentile, from the pooled model and each DA's own
vulnerability profile. [Placeholder---gated on observed
mortality.]}
\label{fig:map}
\end{figure}


%============================================================
\section{Discussion}
%============================================================

Across the dissemination areas of Canada's 41 census
metropolitan areas, heat accounted for an average of [XX] excess
deaths ([95\% eCI]) each year, a standardized excess mortality
rate of [XX] per 100,000 person-years, with minimum-mortality
temperatures ranging from [XX] to [XX]$\,\degree$C across cities.
Risk varied [more within cities than between them], with the
highest values concentrated in [pattern].
% [GATE: all three slots from Results; within/between may flip.]

The patterns align with the small-area literature where it
exists. The component tracking material deprivation carried
[higher] heat risk, consistent with small-area findings in
Europe \citep{gasparrini2022smallarea, zhang2023cvdeurope}. [Density,
green space, and air quality loadings] recur across cities with
only their weights differing \citep{sera2019urban}.
Minimum-mortality temperatures tracked each city's own
distribution, [continental CMAs separating from maritime ones],
the adaptation signature reported across countries
\citep{gasparrini2012mvmeta}. [Renter share], the available
proxy for air conditioning, [behaved as in] the Canadian
multi-city record \citep{quick2024twelvecities}.
% [GATE: every bracket filled from Fig 3 loadings + Results;
%  all refs must appear in the Introduction -- logged for the
%  intro build.]

A map at this resolution places risk at the scale where
heat-health response is organized. Public health units can
identify the most exposed neighbourhoods directly rather than
infer them from city-wide averages, and the fitted risk
functions couple to weather forecasts for anticipatory warning
\citep{mistry2024forecast}. The same framework extends to
hospital admissions and emergency-department visits, and to
projection under climate scenarios. We recommend the
dissemination-area estimates, not city averages, as the input
layer for municipal heat-response planning.

The magnitude of the excess depends on the effect summary. Each
estimate is referenced to the DA's minimum-mortality
temperature, so the excess counts deaths against a counterfactual
at that DA's own optimum; heat and cold burdens are therefore
not comparable across studies using fixed thresholds. The
estimates are averages over the study period and carry no
temporal change in susceptibility.

This analysis has limitations. First, the downscale assumes that
risk patterns estimated across areas hold at the areas
themselves; agreement with unpooled census-tract fits (appendix)
supports the assumption at one scale below the fitting unit.
Second, the composite components size how much variation the
indicators track and cannot attribute it to any single driver.
Third, individual-level factors are not modelled, and the
associations retain an ecological character. In addition,
exposure comes from an interpolated grid whose uncertainty is
not propagated, and predicted curves carry only the between-area
variation the components explain, so the ordering of areas is
more reliable than the range of risk; a surface corrected for
the omitted variancetaken as rates is reported as a sensitivity analysis.


%============================================================
% PANEL: RESEARCH IN CONTEXT (slot even if target isn't Lancet)
%============================================================
% Evidence before this study / Added value / Implications.

%============================================================
% APPARATUS
%============================================================
% \section*{Contributors}
% \section*{Declaration of interests}
% \section*{Acknowledgments}

\bibliographystyle{unsrtnat}
\bibliography{refs}

\end{document}